\documentclass[a4paper,10pt]{article}
\usepackage[T1]{fontenc}
\usepackage{lmodern,microtype,amsmath,amssymb,graphicx,booktabs,tabularx,array}
\usepackage[textwidth=170mm,top=20mm,bottom=20mm]{geometry}
\usepackage[font=small,labelfont=bf]{caption}
\usepackage[colorlinks=true,linkcolor=blue,urlcolor=blue,citecolor=blue]{hyperref}
\pagestyle{plain}\setlength{\parindent}{0pt}\setlength{\parskip}{5pt}
\renewcommand{\arraystretch}{1.15}
\setcounter{tocdepth}{2}
\hypersetup{pdftitle={When Background Noise Has Meaning},pdfauthor={Qihan Zhu}}
\newcommand{\site}{https://charlie-uni.github.io/elec5305-project-550489593/}
\newcommand{\repo}{https://github.com/Charlie-Uni/elec5305-project-550489593}
\newcommand{\pic}[2]{\begin{center}\begin{minipage}{\linewidth}\includegraphics{../results/final_submission_20261110/figures/#1.pdf}\captionof{figure}{#2}\label{fig:#1}\end{minipage}\end{center}}
\newcommand{\CleanRows}{
SVM & 48.90 $\pm$ 4.21 & 47.96 $\pm$ 4.07 \\
CNN & 67.90 $\pm$ 3.65 & 66.61 $\pm$ 3.91 \\
CLAP & 81.75 $\pm$ 2.58 & 79.49 $\pm$ 3.28 \\
}
\newcommand{\AugRows}{
Clean & 65.78 $\pm$ 0.75 & 7.48 $\pm$ 0.42 & 30.92 $\pm$ 1.35 & 32.99 $\pm$ 0.34 & 31.41 $\pm$ 0.51 \\
White noise & 64.65 $\pm$ 0.61 & 43.95 $\pm$ 2.67 & 33.80 $\pm$ 2.19 & 33.46 $\pm$ 0.98 & 32.10 $\pm$ 1.05 \\
Semantic & 61.70 $\pm$ 2.40 & 9.91 $\pm$ 0.24 & 33.62 $\pm$ 1.55 & 36.05 $\pm$ 0.83 & 33.73 $\pm$ 0.86 \\
}
\newcommand{\CostRows}{
SVM & 15.00 & 11.29 \\
CNN & 3.07 & 2.44 \\
CLAP & 69.30 & 50.86 \\
}

\newcommand{\LocalAlignmentRows}{
Single prompt & 81.75 & 52.80 & 48.86 & 632 & 312 & 70.00 \\
Fixed ensemble & 92.90 & 58.30 & 55.50 & 645 & 189 & 78.20 \\
Local: preset & 92.90 & 58.25 & 55.46 & 644 & 191 & 76.90 \\
Local: calibrated & 92.85 & 58.20 & 55.47 & 645 & 191 & 76.10 \\
}

\begin{document}
\fontsize{10.5}{13}\selectfont
\thispagestyle{empty}
\begin{center}
{\LARGE\bfseries When Background Noise Has Meaning}\\[7pt]
{\large Environmental Sound Classification under Semantic Interference}\\[5pt]
ELEC5305 Acoustics, Speech and Signal Processing\\[5pt]
Qihan Zhu\quad SID 550489593
\end{center}
\vspace{3pt}
\begin{abstract}
Environmental recordings often contain multiple recognisable events, yet single-label classifiers must select one. This study compares how three audio representations respond to background noise and competing events at controlled relative power. Across the five official ESC-50 folds, an SVM with handcrafted features, a compact log-Mel CNN and frozen single-prompt LAION-CLAP are evaluated on identical, reproducibly generated audio mixtures. Under the own-activity power convention at +5 dB, white noise reduces SVM macro-F1 to 2.0\%, while CLAP retains 62.9\%; a distractor from the same broad category reduces CLAP to 48.9\%. Across same- and different-category pairings at this level, 60--67\% of CLAP's errors select the distractor. White-noise augmentation transfers weakly to competing events: across three CNN seeds, macro-F1 increases from 7.5\% to 44.0\% under white noise, but only from 33.0\% to 33.5\% under same-category interference. Local audio--text feedback on the same frozen checkpoint provides no statistically significant accuracy gain over fixed prompt ensembling in the primary comparisons. The target remains second in approximately one-third of cases in which CLAP selects the distractor, showing that some first-place errors retain the target among the leading candidates. In small subsets matched on measured acoustic descriptors, CLAP remains less accurate under competing events than under DEMAND backgrounds, while the corresponding CNN contrasts are inconclusive. These findings indicate that noise robustness does not guarantee target retention and that candidate rankings help interpret classification errors. Synthetic mixtures, single-label evaluation, limited pairings and backgrounds, and restricted acoustic matching constrain generalisation; the results do not isolate the causal contribution of semantic content.
\end{abstract}
\vspace{3pt}
{\small GitHub Username: Charlie-Uni\\[1pt]
GitHub Project Link: \url{https://github.com/Charlie-Uni/elec5305-project-550489593}\\[1pt]
Web: \href{\site}{Project website}\quad Video: \href{https://charlie-uni.github.io/elec5305-project-550489593/results/final_submission_20261110/project_video.mp4}{Project video report}\par}
\vspace{3pt}
\begingroup\small\setlength{\parskip}{0pt}\tableofcontents\endgroup
\clearpage
\section{Introduction}

Single-label environmental sound classification requires one prediction per recording, yet real acoustic environments often contain overlapping events \cite{sed}. Interference may be acoustically nonspecific, such as white noise, or another recognisable event. For example, with dog barking as the designated target, white noise adds acoustic interference, whereas a ringing alarm introduces another plausible event label. The two mixtures can have the same power ratio yet produce different errors. Work on structured acoustic scenes motivates looking beyond Gaussian noise \cite{rsa}. Spectrogram CNNs and audio--language pretraining offer different representations \cite{piczak,clap}, but benchmark performance alone does not establish how predictions shift towards a known competing class under identical mixtures. This study uses all 2,000 five-second ESC-50 clips, covering 50 classes assembled from Freesound \cite{esc}, as targets, draws competing events from this benchmark subject to the source exclusions in Section~\ref{sec:data}, and uses DEMAND recordings \cite{demand} as real backgrounds alongside generated white noise.

The main research question is: \emph{How do handcrafted, supervised spectrogram and pretrained audio--language representations fail under generic noise versus a competing labelled event at controlled relative power?} Error destinations and the transfer of noise augmentation are subordinate tests of this question. The intended contribution is to explain measured failure patterns rather than identify a universally best model. In particular, CLAP's large external pretraining makes equal-data or equal-training-cost comparisons inappropriate.

Every system receives the same mixtures. Each mixture has a clean control with the same gain, allowing changes induced by interference to be distinguished from errors already present after rescaling. The study covers one benchmark, three fixed systems, four interference types and three relative levels. The proposal's target-activity power convention is retained as a sensitivity analysis.

\section{Literature review}

\subsection{Course foundations and related research}

Sound event detection identifies which events occur and when they are active, including overlapping events \cite{sed}, whereas single-label classification assigns one class to an entire recording. The course material introduces energy, zero-crossing rate, spectral descriptors and MFCCs as interpretable features, although temporal aggregation discards event order and may combine several sources \cite{course3}. The STFT frontend further determines the available information: window length, hop size and FFT length govern the time--frequency trade-off, frame spacing and frequency-grid spacing, respectively. Zero padding refines the grid without recovering detail unresolved by the window \cite{course4}.

Piczak's CNN learns local patterns from spectrogram segments and their deltas, complementing fixed descriptors with trainable features \cite{piczak}. Its supervision remains limited by the amount of labelled data. Historical ESC-50 accuracies of 39.60\% for MFCC/ZCR SVM and 64.50\% for Piczak's CNN provide context rather than matched comparisons \cite{piczak,esc}. Augmentation broadens training conditions, but its benefits depend on the transformation and class \cite{augmentation}. BC learning mixes different-class sounds and learns their proportions \cite{bc}, while mixup generalises interpolation of inputs and labels \cite{mixup}. Both construct soft labels from mixing proportions; this study retains the designated target label to test recognition under interference.

Audio--language pretraining changes the source of supervision. CLAP matches audio with class descriptions without fitting an ESC-50 classifier \cite{clap}, although pretraining data and prompt wording influence performance. Its audio--text similarity also provides a way to examine competition beyond the final prediction: target and distractor scores and rankings can be tracked as interference strength changes. PAT further improves cross-modal alignment through prompt weighting and frame-level audio feedback, with evaluations across several audio--language models and on degraded audio \cite{pat}.

Recent work extends robustness evaluation beyond clean benchmarks. Mix-CLAP adaptively fuses embeddings from encoders distilled on clean and noisy audio, addressing robustness alongside deployment efficiency \cite{mixclap}. RSA-Bench studies speech perception and reasoning under structured, multisource backgrounds \cite{rsa}. AIP-Speech relates acoustic descriptors, including zero-crossing rate, spectral centroid and roll-off, to degradation in speech recognition and keyword spotting \cite{aip}. Because RSA-Bench and AIP-Speech evaluate speech tasks, they support the experimental rationale rather than serve as directly comparable ESC-50 classification baselines.

\subsection{Evidence gaps and methodological choices}

Three gaps in the existing evidence motivate the experimental controls. Augmentation benefits depend on the transformation and class \cite{augmentation}, while BC learning and mixup reward predictions that reflect both mixed labels \cite{bc,mixup}. Those objectives support regularisation and recognition of mixtures, but do not directly measure retention of one designated event. This study therefore compares white-noise and event-mixture augmentation using the same target-label objective, pairing rules and test mixtures. Gain-matched clean controls distinguish new errors from those already present after rescaling.

Audio--language alignment presents a second distinction: matching a recording to plausible descriptions does not specify which of two present events should rank first. PAT's feedback improves cross-modal agreement \cite{pat}, but preservation of a designated target against a valid competing label requires a separate test. Holding the LAION checkpoint frozen across single-prompt, fixed-ensemble and feedback readouts makes their effects comparable. The inspected PAT release uses MS-CLAP 2023 \cite{pat,msclap}; this project adapts its arithmetic to LAION with uniform prompt weights and project-selected projection and coefficient rules (Section~\ref{sec:metrics}). Calibration uses non-test clean labels and is reported separately from zero-shot readouts.

Structured-background evaluations broaden acoustic coverage \cite{rsa,aip}, but speech-task scores do not resolve ESC-50 error destinations. Characterising a background independently also leaves its overlap with a particular target unmeasured. Pairwise spectral and temporal descriptors, followed by same-target acoustic matching, address that narrower comparison. Project code implements STFT, features, the compact CNN and mixing and evaluation using external numerical and learning libraries. These inspectable representations complement frozen CLAP, while separate CPU timings retain the deployment concern raised by Mix-CLAP \cite{mixclap}.

These controls address measured acoustics, not perceptual equivalence: frequency-dependent hearing allows equal power ratios to yield different prominence \cite{course6}. Class scores and rankings describe relative matches rather than calibrated event probabilities. Regressions with and without acoustic covariates examine associations with classification success; non-collapsibility also affects coefficient changes, which are interpreted descriptively \cite{noncollapsibility}.

\section{Methods}

\subsection{Data and evaluation splits}\label{sec:data}

ESC-50 contains 2,000 five-second clips, 50 balanced classes and five official folds \cite{esc}. For test fold $k$, the next cyclic fold is validation and the remaining three support model selection. SVM and CNN models are then refitted on the four non-test folds. Scaling and Mel normalisation use the corresponding training partition. Four source IDs cross folds, affecting nine target clips; official assignments are retained, but these sources are excluded as semantic distractors. Residual source overlap and partly class-dependent preprocessing remain dataset limitations \cite{esc}.

\subsection{Signal processing and classifiers}

At 44.1 kHz, the custom STFT uses a periodic Hann window of 1,024 samples, hop 512 and FFT length 1,024:

\[
X(m,k)=\sum_{n=0}^{1023}x[n+512m]w[n]e^{-j2\pi kn/1024}.
\]

Only the final incomplete frame is zero-padded. Comparison with SciPy and an actual execution of the course's unchanged MyStft.m gives a relative complex error of approximately $2.6\times10^{-16}$ on one ESC-50 clip after matching frame alignment, windowing and one-sided output. This comparison checks corresponding frames after harmonising the implementations; their default boundary policies can differ.

\begin{center}\begin{minipage}{\linewidth}\small\centering
\captionof{table}{Classifier representations, fitting and scoring. Supervised hyperparameters are selected using clean validation data.}\label{tab:classifiers}

\begin{tabularx}{\textwidth}{@{}lXX@{}}\toprule
System & Representation & Fitting and scoring\\\midrule
SVM & Means/SDs of RMS, ZCR, centroid, spread, roll-off, flux, 20 MFCCs and deltas: 92 values & RBF kernel; training-only scaling; $C\in\{1,10,100\}$; $\gamma$: \texttt{scale}, 0.001 or 0.1; nine combinations selected by clean validation\\
CNN & 64 log-Mel bands; three convolution blocks (16/32/64 channels); 256-dimensional embedding & 36,258 parameters; Adam; at most 40 epochs, patience 8; clean validation chooses duration\\
CLAP & HTSAT-tiny/RoBERTa, non-fusion \texttt{630k-audioset-best.pt}; unit embeddings & Frozen; 48-kHz resampling; official int16 truncation; cosine scores for \texttt{the sound of \{label\}}\\\bottomrule
\end{tabularx}
\end{minipage}\end{center}

Classifier configurations are summarised in Table~\ref{tab:classifiers}. CLAP repeat-pads five-second inputs to ten seconds and uses one fixed prompt template across classes. The core CNN seed is 5305; repetitions use 5306 and 5307. Scores and representation distances are interpreted within each model.

\subsection{Controlled mixing and power sensitivity}

Every target receives white noise, a segment from one of three DEMAND recordings (traffic, station or metro) \cite{demand}, and one same- and one different-coarse-category ESC-50 distractor. Semantic partners differ in fine class and source and come from the held-out fold. Seeded pairing uses zero offset and remains fixed across +10, +5 and 0 dB; noise waveforms and segments also remain unchanged.

Let $x$ and $d$ denote target and interferer, $M_x$ and $M_d$ their activity masks, and $P(z;M)$ mean-square power over mask $M$. Active frames lie within 40 dB of each signal's peak frame energy. Active-region power reduces the underestimation of sparse events, such as knocks, that can arise from averaging over the full recording. $L$ denotes SNR for noise and SIR for labelled-event mixtures:

\[
\alpha=\sqrt{\frac{P(x;M_x)}{10^{L/10}P(d;M)}},\qquad y=g(x+\alpha d),\qquad g=\min\left(1,\frac{0.99}{\max_n|x[n]+\alpha d[n]|}\right).
\]

The factor $\alpha$ sets the interference level, while $g$ attenuates both sources together to limit the mixture peak to 0.99. The control is $gx$. The default own protocol uses $M=M_d$; the proposal's target protocol uses $M=M_x$. Large gains for sparse distractors motivated the documented change and full comparison of both protocols. Achieved ratios are verified numerically; manifests retain source labels, seeds, offsets, mixing gains and common attenuation. Own contains 24,000 mixtures and target 23,991: three pairs with zero target-region interference power exclude nine target rows. Equal-power 0-dB mixtures serve as stress tests because a single designated label is ambiguous.

\subsection{Metrics, uncertainty and extension design}\label{sec:metrics}

Accuracy and macro-F1 are averaged over official folds; macro-F1 reflects performance across classes when predictions concentrate on a few labels. Capture rate is the fraction of errors assigned to the designated distractor; capture counts across all trials are also retained. Margin is target score minus distractor score, and top-2 analysis records target and distractor inclusion. Representation displacement compares mixtures with gain-matched controls, using training-fold feature standardisation or within-model embedding cosine distance. Logistic regressions with and without acoustic covariates describe associations between interference type and classification success; coefficient changes are interpreted descriptively in light of non-collapsibility \cite{noncollapsibility}.

Paired comparisons reuse target clips. Ranking, acoustic-matching and local-feedback analyses use 2,000 bootstrap draws clustered by original target recording, with pointwise 95\% intervals conditional on fitted models, calibration rules and observed pairings; shared interferers remain a dependency. Earlier high-overlap replication uses clip-level paired intervals. Fold and seed SDs describe variability rather than confidence intervals. Holm correction covers separate families of 12 replication and 156 intervention tests; the latter was revised after early white-noise results but before semantic results.

SVM/CNN augmentation draws only from training partitions. CNN augmentation probability is 0.5 at +5/+10/+15 dB; architecture and clean validation remain fixed, with all three seeds retained. A fixed ensemble of 300 prompts reuses CLAP audio embeddings without test-label prompt selection or fine-tuning.

Local feedback follows the released PAT arithmetic: 32 frame features pass through the model's audio projection, and bidirectional text--frame attention produces text- and audio-guided contributions added to the prompt-ensemble score. The readout is reported as a fused score, distinct from cosine similarity. Preset temperature 64 and text-/audio-guided weights 0.2/0.05 are heuristic choices fixed before evaluation without labels. For each test fold, calibration selects among 321 settings using other folds' clean clips, excluding shared source recordings; one setting applies to every interference condition in that fold. This variant is a label-calibrated readout and is reported separately from zero-shot results. Four primary comparisons share one Holm correction: both variants against fixed ensembling under same- and different-category interference at own/+5 dB.

\section{Results}

\subsection{Core robustness and error destinations}

\begin{center}\begin{minipage}{\linewidth}\small\centering
\captionof{table}{Clean performance: fold mean $\pm$ SD (\%). Capture: pooled captures/errors at own/+5 dB, with 2,000 targets per condition, CNN seed 5305 and single-prompt CLAP. The capture-rate denominator is errors, not all trials.}\label{tab:results2}
\begin{tabular}{@{}lrrrr@{}}\toprule
 & \multicolumn{2}{c}{Clean (\%)} & \multicolumn{2}{c}{Capture (\%; captures/errors)}\\
System & Accuracy & Macro-F1 & Same category & Different category\\\midrule
\CleanRows\bottomrule\end{tabular}
\end{minipage}\end{center}

\pic{01_robustness}{Own-activity robustness: mean macro-F1 with five-fold SD, not a confidence interval. Each condition and level contains 2,000 targets. Lower relative levels indicate stronger interference. Target-protocol curves are available in the repository.}

At +5 dB (Figure~\ref{fig:01_robustness}), SVM macro-F1 is 2.03\% under white noise versus 22.47\% under same-category interference. CLAP instead reaches 71.49\% under DEMAND versus 48.86\% under same-category interference. Equal nominal power therefore produces different degradation across interference types.

Error destinations distinguish the systems (Table~\ref{tab:results2}). The designated distractor receives 66.95\% and 60.30\% of CLAP's errors, against 18.53\% and 15.89\% for the SVM; uniform confusion over 49 wrong labels would give about 2\%. For CLAP at own/+5 dB, the same-category median target--distractor cosine margin is 0.0873, with 38.05\% negative, and median embedding displacement is 0.1905 for correct versus 0.4087 for incorrect trials. For the SVM, median standardised ZCR displacement is 3.16 under white noise against 0.25--0.29 under the other three conditions, consistent with broadband noise changing zero crossings and spectral-shape descriptors.

Figure~\ref{fig:04_capture_margin} extends these comparisons across levels. Capture rises as SIR falls for every system, from 56.25\% to 71.45\% of CLAP's same-category errors, while median margins move towards zero at 0 dB. The most frequent CLAP captures at +5 dB among pairs whose control was correct, selected from observed counts, are door knock to keyboard typing (10 of 11 pairs) and can opening to clock alarm (6 of 6). Figure~\ref{fig:05_displacement} shows larger median displacement among misclassified mixtures in all 12 system/condition comparisons, although the individual distributions overlap. White noise displaces the SVM and CNN most, whereas CLAP's misclassified semantic mixtures move about as far as its misclassified white-noise mixtures, with medians of 0.41--0.43.

\pic{04_capture_margin}{Own-activity semantic interference across levels. (a) Share of each system's errors assigned to the designated distractor; the dotted line represents uniform confusion over 49 wrong labels. (b--d) Median target-minus-distractor score with interquartile range, using a separate scale for each system. Filled/open markers indicate same/different category; panel (a) also uses solid/dashed lines. Each level contains 2,000 targets, and margin bars show the score IQR.}

\pic{05_displacement}{Displacement of each mixture from its gain-matched control at own/+5 dB: median and interquartile range for correctly classified (filled) and misclassified (open) mixtures. Each condition contains 2,000 trials partitioned by outcome: correct groups contain 108--1,479 trials and error groups 521--1,892. These are descriptive, outcome-conditioned distributions, with separate scales across systems.}

\subsection{Overlap, training interventions and readout}

Across all targets at own/+5 dB, the core CNN has higher macro-F1 under a same-category event than under DEMAND (32.60\% versus 29.46\%), but many pairs overlap little in time. When at least 80\% of target-active samples also contain active interference, paired semantic-minus-DEMAND accuracy differences are negative for all three seeds. For seed 5305, the differences are $-3.28$ percentage points on 1,464 same-category targets and $-2.94$ points on 1,465 different-category targets, with Holm-adjusted $p=0.0604$ for both comparisons. The direction replicates in all three seeds, but only the two additional seeds pass Holm correction. This subset comparison uses accuracy rather than the whole-set macro-F1 above. Target-power semantic macro-F1 is 0.47--2.35 points lower than own power; slightly different retained cohorts prevent attribution solely to the power mask.

\begin{center}\begin{minipage}{\linewidth}\small\centering
\captionof{table}{CNN augmentation: mean macro-F1 $\pm$ sample SD across three seeds. Own-power interference conditions are evaluated at +5 dB; clean evaluation uses the original clips.}\label{tab:results3}
\begin{tabular}{@{}lrrrrr@{}}\toprule Training & Clean & White & DEMAND & Same & Different\\\midrule\AugRows\bottomrule\end{tabular}
\end{minipage}\end{center}

White-noise augmentation changes CNN white-noise macro-F1 (Table~\ref{tab:results3}) from 7.48\% to 43.95\%, but same-category performance only from 32.99\% to 33.46\%. Semantic training reaches 36.05\% under same-category interference while clean performance falls to 61.70\%. The direct accuracy advantages of semantic over white-noise augmentation do not pass the 156-test correction. SVM augmentation results and all individual seeds are retained in the repository.

Fixed prompt ensembling raises CLAP clean accuracy from 81.75\% to 92.90\%. At own/+5 dB, same-category captures change from 632 to 645 while other errors fall from 312 to 189. The net improvement therefore comes from fewer other-class errors, not fewer captures; a larger fraction of the remaining errors can name the distractor despite higher accuracy.

\begin{center}\begin{minipage}{\linewidth}\small\centering
\captionof{table}{CLAP readouts on clean audio and own/+5-dB same-category interference, each with 2,000 trials. Accuracy, macro-F1 and target top-two inclusion are percentages; captures and other errors are counts. The calibrated readout uses non-test clean labels.}\label{tab:results4}
\begin{tabular}{@{}lrrrrrr@{}}\toprule Readout & Clean acc. & Same acc. & Same F1 & Captures & Other errors & Target top-2\\\midrule\LocalAlignmentRows\bottomrule\end{tabular}
\end{minipage}\end{center}

The readouts in Table~\ref{tab:results4} use a new extraction path that exactly reproduces the original single-prompt and fixed-ensemble predictions. For the local-feedback port in Section~\ref{sec:metrics}, preset feedback changes accuracy relative to the ensemble by $-0.05$ percentage points for same-category and +0.20 points for different-category interference at own/+5 dB; calibration changes it by $-0.10$ and +0.75 points, respectively. None of the four primary tests passes Holm correction (smallest adjusted $p=0.735$), and all pointwise source-cluster intervals include zero. Same-category captures remain at 644--645. Across all 2,000 same-category trials, target top-two inclusion is 78.20\% for the ensemble, 76.90\% for preset feedback and 76.10\% after calibration. Thus, no improvement in target retention is confirmed for this port.

\pic{02_top_two}{Single-template CLAP at own/+5 dB. Exact top-two membership of the two designated labels among all 2,000 semantic trials per type. The four outcomes partition the trials; returned extra labels are not exhaustively annotated events.}

Among single-prompt captures, the target remains second in 236/632 (37.34\%) same-category and 203/612 (33.17\%) different-category cases. These denominators differ from the all-trial percentages in Figure~\ref{fig:02_top_two}. Gain-matched controls contain both labels in the top two in only 5.40\% and 1.70\%, respectively. Target-power results show similar partial recovery.

\subsection{Spectral and temporal overlap}

Spectral and temporal descriptors cover all 47,991 mixtures, including 12,000 DEMAND rows. The repository retains additional descriptor checks and sensitivity analyses.

To compare the same target under semantic interference and DEMAND, the primary analysis retains pairs with differences of at most 0.10 in target-active spectral cosine, 0.05 in the interference-dominated time--frequency fraction and 0.10 in activity overlap. Selection uses acoustics only, identically for all three systems. These are acoustic proxies from the project frontend, not CLAP's internal frontend or a perceptual masking model.

\pic{03_acoustic_matching}{Semantic-minus-DEMAND accuracy at +5 dB, with pointwise 95\% target-source cluster intervals from 2,000 resamples. Own: 210/178 targets, 199/168 source clusters and 49/48 classes; target: 205/187 targets, 195/178 clusters and 49/47 classes, listed as same/different category. Filled circles indicate same category; open squares indicate different category.}

Across both power conventions (Figure~\ref{fig:03_acoustic_matching}), CLAP remains about 16--26 accuracy percentage points below DEMAND, whereas all four CNN intervals include zero. Only about 9--10.5\% of the benchmark remains. These selected populations differ from the larger high-overlap subsets; the comparison is not a causal decomposition or a reversal on identical targets.

\subsection{Deployment cost}

\begin{center}\begin{minipage}{\linewidth}\small\centering
\captionof{table}{Median CPU query time (ms) on the same Apple M3, using one numerical thread and batch size one. Each pass contains 250 class/fold-balanced clips and three repetitions per model. Both model orders are retained.}\label{tab:results5}
\begin{tabular}{@{}lrr@{}}\toprule System & First pass & Reversed order\\\midrule\CostRows\bottomrule\end{tabular}
\end{minipage}\end{center}

Query cost (Table~\ref{tab:results5}) includes preprocessing and scoring, excluding file access, model loading and one-off preparation. Absolute times vary, but the CNN/SVM/CLAP ordering persists. CNN and CLAP have 36,258 and 158,334,169 parameters, respectively; external pretraining is excluded from the cost measurements. Local feedback adds about 2.6--2.7 ms per query in a separate identical-clip panel: ensemble medians are 59--60 ms, versus 62--63 ms for either aligned version. Both orders, each with 750 queries per readout, are retained. Encoding clean clips and calibrating on cached scores are recorded as separate one-off costs.

\section{Discussion and conclusion}

\subsection{Interpretation of results}

The central distinction is between degraded acoustic evidence and a decision favouring another event that is actually present. The observed changes in the SVM's zero-crossing and spectral statistics under white noise are consistent with the course descriptors \cite{course3}; combining displacement with score margins shows why movement relative to class boundaries matters alongside its magnitude. The clean SVM and CNN accuracies are 48.90\% and 67.90\%, compared with historical values of 39.60\% and 64.50\%, respectively \cite{piczak,esc}. These provide baseline context under different feature and training configurations. The contribution is the controlled comparison of their subsequent failures.

Augmentation exposes a mismatch between training coverage and deployment demands. Learning invariance to broadband corruption gives limited protection when another event carries class-relevant structure. Semantic augmentation changes that training exposure, but trades some clean performance for semantic point-estimate gains; its direct accuracy advantages over white-noise augmentation do not pass the 156-test correction. This extends the transformation- and class-dependent augmentation evidence \cite{augmentation}, using target retention as a different objective from the mixed-label learning in BC learning and mixup \cite{bc,mixup}.

CLAP's class-description scores \cite{clap} make a competing event a plausible answer. The classifier receives the mixture without an explicit cue identifying the experiment's designated target, so capture measures disagreement with that target-selection rule. Top-2 analysis shows that some first-place changes retain the target among leading candidates. This distinguishes ranking failure from complete exclusion of the target, while leaving the interpretation of perception and event detection to separately validated tasks.

Fixed prompt ensembling primarily improved performance by reducing other-class errors, whereas distractor captures did not fall. Overall accuracy and resistance to capture are therefore different optimisation objectives. PAT's reported LAION-CLAP improvement from 89.00\% to 93.00\% under its own ESC-50 protocol supports the value of alignment \cite{pat}, but does not determine target preference in these mixtures. The local-feedback adaptation adds no confirmed primary gain over fixed ensembling. Its clean-label calibration rewards clean accuracy; target retention against a competing event is a different selection objective, suggesting a testable source of the mismatch.

Temporal stratification and acoustic matching condition the comparison on different populations. High-overlap CNN contrasts share a direction across seeds, whereas intervals in the smaller acoustically matched subsets include zero. This does not show that acoustic adjustment removed an effect on identical targets. CLAP's remaining matched-subset deficit supports a narrower behavioural observation under closer measured acoustics. Pairwise descriptors complement AIP-Speech's background-level associations \cite{aip} and RSA-Bench's structured-scene evaluation \cite{rsa}, while residual differences still require independent acoustic controls for a semantic explanation.

Practically, evaluation should test ordinary noise and competing events separately, with fixed pairings and power conventions. Augmentation and readout selection should use metrics that reflect the intended task, reporting clean performance, target retention and capture together. Where multiple labels are useful, ranked candidates provide a starting point for a calibrated readout whose false positives must also be assessed. Deployment choices should then balance this condition-specific evidence against measured latency: the compact CNN is fastest on the tested hardware, connecting robustness assessment with Mix-CLAP's efficiency concern \cite{mixclap}.

\subsection{Limitations and further validation}

Each target currently has one distractor per category type, and environmental backgrounds come from three DEMAND recordings. Repetition with new pairing seeds and background segments would assess stability across mixture sampling. Recording-start alignment and audible-event overlap are separate conditions; finer offset and SIR sweeps within fixed pairs would help locate decision transitions.

Reuse of interferer recordings across trials, limited source overlap across official folds and uncertain overlap with CLAP's pretraining data limit generalisation. Retaining official folds as the primary evaluation while adding a source-independent sensitivity study would help distinguish these influences. Acoustic matching currently retains approximately 9--10.5\% of targets. Expanding the matched population and independently validating the relationship between matching descriptors and perceptual prominence would support further investigation of acoustic structure and class identity.

The designated target is defined by the single-label experimental protocol, while a multi-event scene contains several plausible labels. Multi-event annotation and calibrated readouts would allow target retention, event detection and temporal activity to be evaluated separately \cite{sed}. For local feedback, mixed calibration examples from non-test folds could support controlled comparisons of prompt weighting, frame projection and calibration objectives. Future repetitions should prespecify power conventions, top-2 metrics and matching thresholds, turning choices developed from early results into a prospective validation protocol.

\subsection{Conclusion}

Using the official ESC-50 folds and identical mixtures, this project compares how three audio representations fail under ordinary noise and competing events. At own/+5 dB, white noise reduces SVM macro-F1 to 2.03\%, while single-prompt CLAP retains 62.91\%; CLAP scores 48.86\% under same-category interference versus 71.49\% under DEMAND. Across both semantic conditions, 60--67\% of its errors select the distractor, while approximately one-third of single-prompt captures retain the target at rank two. Averaged across three CNN seeds, white-noise augmentation raises white-noise macro-F1 from 7.48\% to 43.95\%, but same-category macro-F1 only from 32.99\% to 33.46\%, indicating limited transfer to competing events. Fixed prompt ensembling raises clean CLAP accuracy to 92.90\%, while same-category captures change from 632 to 645; local feedback provides no confirmed target-retention gain over fixed ensembling in the primary comparisons.

The project contributes a reproducible experiment connecting STFT, audio features, power control and classification diagnostics. Within the synthetic-mixture and single-label protocol studied here, deployment evaluation should distinguish noise robustness from target retention under competing events and use candidate rankings to interpret first-place errors. The extent to which CLAP's competing-class preference reflects semantic content remains a question for further investigation using independently validated acoustic controls.

\section*{AI acknowledgement}
\addcontentsline{toc}{section}{AI acknowledgement}
I used OpenAI Codex (GPT-6.1 Sol; \url{https://openai.com/codex/}) and Anthropic Claude (Fable 5.1; \url{https://claude.ai/}) to assist with parts of the code implementation and revision, experiment orchestration, numerical checking, figure and document production, and drafting. I was responsible for the project's overall conception and requirements analysis, selecting and approving the methods, reviewing the key implementation and reported results, and developing the interpretation and argument presented in the final submission.

\clearpage\phantomsection\addcontentsline{toc}{section}{References}
\begin{thebibliography}{17}\raggedright
\bibitem{sed} A. Mesaros, T. Heittola, T. Virtanen and M. D. Plumbley, ``Sound Event Detection: A Tutorial,'' IEEE Signal Processing Magazine, vol. 38, no. 5, pp. 67--83, 2021. \href{https://doi.org/10.1109/MSP.2021.3090678}{doi:10.1109/MSP.2021.3090678}.
\bibitem{rsa} Y. Zhang et al., ``RSA-Bench: Benchmarking Audio Large Models in Real-World Acoustic Scenarios,'' Findings of ACL, 2026, pp. 18350--18374. \url{https://aclanthology.org/2026.findings-acl.913/}.
\bibitem{piczak} K. J. Piczak, ``Environmental Sound Classification with Convolutional Neural Networks,'' MLSP, 2015. \href{https://doi.org/10.1109/MLSP.2015.7324337}{doi:10.1109/MLSP.2015.7324337}.
\bibitem{clap} Y. Wu et al., ``Large-Scale Contrastive Language-Audio Pretraining with Feature Fusion and Keyword-to-Caption Augmentation,'' ICASSP, 2023. \href{https://doi.org/10.1109/ICASSP49357.2023.10095969}{doi:10.1109/ICASSP49357.2023.10095969}; revised \href{https://arxiv.org/abs/2211.06687v4}{arXiv v4}, 2024.
\bibitem{esc} K. J. Piczak, ``ESC: Dataset for Environmental Sound Classification,'' ACM Multimedia, 2015. \href{https://doi.org/10.1145/2733373.2806390}{doi:10.1145/2733373.2806390}. Official folds and caveats: \url{https://github.com/karolpiczak/ESC-50}.
\bibitem{demand} J. Thiemann et al., ``The Diverse Environments Multi-Channel Acoustic Noise Database (DEMAND),'' POMA, 2013. \href{https://doi.org/10.1121/1.4799597}{doi:10.1121/1.4799597}; data: \href{https://doi.org/10.5281/zenodo.1227121}{Zenodo 1227121}.
\bibitem{course3} C. Jin, \emph{ELEC5305 Lecture 3: Sound Synthesis Models and Audio Features}, University of Sydney, 2026. Course material, pp. 33--38, 50--51, 62--64.
\bibitem{course4} C. Jin, \emph{ELEC5305 Lecture 4: STFT}, University of Sydney, 2026. Course material, pp. 2, 15, 24--25, 31, 87; accompanying \texttt{MyStft.m}.
\bibitem{augmentation} J. Salamon and J. P. Bello, ``Deep Convolutional Neural Networks and Data Augmentation for Environmental Sound Classification,'' IEEE Signal Processing Letters, 2017. \href{https://doi.org/10.1109/LSP.2017.2657381}{doi:10.1109/LSP.2017.2657381}.
\bibitem{bc} Y. Tokozume, Y. Ushiku and T. Harada, ``Learning from Between-class Examples for Deep Sound Recognition,'' ICLR, 2018. \url{https://arxiv.org/abs/1711.10282}.
\bibitem{mixup} H. Zhang, M. Cisse, Y. N. Dauphin and D. Lopez-Paz, ``mixup: Beyond Empirical Risk Minimization,'' ICLR, 2018. \href{https://openreview.net/forum?id=r1Ddp1-Rb}{Conference record}; preprint \href{https://doi.org/10.48550/arXiv.1710.09412}{doi:10.48550/arXiv.1710.09412}.
\bibitem{pat} A. Seth et al., ``PAT: Parameter-Free Audio-Text Aligner to Boost Zero-Shot Audio Classification,'' NAACL, 2025, pp. 12376--12394. \url{https://aclanthology.org/2025.naacl-long.616/}.
\bibitem{mixclap} W. Kohno et al., ``Mix-CLAP: Adaptive Fusion of Knowledge-Distilled Audio Embeddings for Noise-Aware Audio-Language Models,'' ICASSP, 2026, pp. 20501--20505. \href{https://doi.org/10.1109/ICASSP55912.2026.11462496}{doi:10.1109/ICASSP55912.2026.11462496}.
\bibitem{aip} A. Dsouza et al., ``AIP-Speech: Auditing Speech LLM Robustness to Background Interference,'' IMPACT-SPEECH Workshop, pp. 10--19, 2026. \href{https://doi.org/10.18653/2026.impactspeech-1.2}{doi:10.18653/2026.impactspeech-1.2}.\bibitem{msclap} B. Elizalde, S. Deshmukh and H. Wang, ``Natural Language Supervision for General-Purpose Audio Representations,'' ICASSP, 2024, pp. 336--340 (MS-CLAP 2023). \href{https://doi.org/10.1109/ICASSP48485.2024.10448504}{doi:10.1109/ICASSP48485.2024.10448504}.
\bibitem{course6} C. Jin, \emph{ELEC5305 Lecture 6}, University of Sydney, 2026. Course material, pp. 2--5: hearing sensitivity and frequency weighting.
\bibitem{noncollapsibility} N. A. Schuster et al., ``Noncollapsibility and its role in quantifying confounding bias in logistic regression,'' BMC Medical Research Methodology, 2021. \href{https://doi.org/10.1186/s12874-021-01316-8}{doi:10.1186/s12874-021-01316-8}.
\end{thebibliography}
\end{document}
